
\documentclass[a4paper,12pt]{letter}
\usepackage[utf8]{inputenc}
\usepackage[english,slovene]{babel}
%\usepackage{blindtext}
%\usepackage{pdflscape}
\usepackage[tikz]{bclogo}
\usepackage{qrcode,marvosym}
\usepackage{pgf,tikz,pgfplots}
\pgfplotsset{compat=1.14}
\usepackage{mathrsfs}
\usetikzlibrary{arrows}
\usepackage{graphicx}
\newcommand{\ds}{\displaystyle}
\usepackage{fancybox}
\usepackage{amsfonts}
\usepackage{amsmath}
%\usepackage{color}
\usepackage{wallpaper}
\usepackage{hyperref}
\usepackage[cp1250]{inputenc}
\usepackage{array}%%%%%%%%%%%%%%%%%%%%%%%%%%%%%%%%%%%%%%%%%%%%
\newcolumntype{M}[1]{>{\centering\arraybackslash}m{#1}}%%%%
\newcolumntype{N}{@{}m{0pt}@{}}%%%%%%%%%%%%%%%%%%%%%%%%%%%%
   \usepackage{fancyhdr,enumerate}

 \newcommand{\ora}{\overrightarrow}
 \renewcommand{\headrulewidth}{3pt}
 \renewcommand{\headrule}{\hbox to\headwidth{%
   \color{gray}\leaders\hrule height \headrulewidth\hfill}}
    \renewcommand{\footrulewidth}{2pt}
    \renewcommand{\footrule}{\hbox to\headwidth{%
  \color{brown}\leaders\hrule height \footrulewidth\hfill}}
\usepackage{gensymb}
\usepackage{amssymb} 

  \usepackage{fancyhdr}
  \textwidth 20 true cm
  \oddsidemargin -2 true cm
  \evensidemargin -2 true cm
  \topmargin -3.5 true cm
  \textheight 27.5 true cm
  \linespread{1.2}
    \renewcommand{\headrulewidth}{0.3pt}
   \renewcommand{\footrulewidth}{1.9pt}
  \definecolor{qqwuqq}{rgb}{1,0,0 }
  \definecolor{qqqqff}{rgb}{0,0,1}
  \definecolor{qqffqq}{rgb}{0,1,0}
  \definecolor{ffqqqq}{rgb}{1,0,0}
\definecolor{zzuxux}{rgb}{0,0,0}%{0.92,0.03,0.03}
\usepackage{mdframed}
\usepackage{multicol}
 \definecolor{bgblue}{RGB}{0,0,253}
\usepackage{xcolor}
\usepackage{eurosym}
%%%%%%%%%%%%%%%%%%%%%%%%%%%%%%%%%%%%%%%%%%%%%%%%%%%%%%%%
\newcounter{tocke}
\setcounter{tocke}{0}
\usepackage{tikz-3dplot}
%%%%%%%%%%%%%%%%%%%%%%%%%%%%%%%%%%%%%%%%%%%%%%%%%%%%%%%%%%
\mdfdefinestyle{theoremstyle}{%
linecolor=brown!40,linewidth=1pt,%
frametitlerule=true,%
frametitlebackgroundcolor=brown!50,
innertopmargin=\topskip,
}

\mdfdefinestyle{theoremstyle1}{%
linecolor=brown,middlelinewidth=1pt,%
frametitlerule=true,%
apptotikzsetting={\tikzset{mdfframetitlebackground/.append style={%
shade,left color=black!40, right color=gray!10},
mdfbackground/.append style={
top color=white!100!white,
bottom color=black!5
},
}}, 
frametitlerulecolor=gray,
frametitlerulewidth=2pt,
innertopmargin=\topskip,
innerbottommargin=\topskip,
}
\mdtheorem[style=theoremstyle1]{naloga}{Naloga}
\mdtheorem[style=theoremstyle]{resitev}{Analiza Naloge}
\mdtheorem[style=theoremstyle]{kriterij}{\v{S}tevilo dose\v{z}enih to\v{c}k na testu}
%%%%%%%%%%%%%%%%%%%%%%%%%%%%%%%%%%%%%%%%%%%%%%%%%%%%%%%%%%
\usepackage[table]{xcolor}
\usepackage{venndiagram}
\usepackage[printwatermark]{xwatermark}
\usepackage{xcolor}
\newwatermark[allpages,color=brown!2,angle=45,scale=5,
xpos=-5,ypos=0]{matej.info}
%%%%%%%%%%%%%%%%%%%%%%%%%%%%%%%%%%%%%%%%%%%%%%%%%%%%%%%%%%
%%%%%%%%%%%%%%%%%%%%%%%%%%%%%%%%%%%%%%%%%%%%%%%%%%%%%%%%%%
\begin{document}
\pagestyle{fancy}
\renewcommand\headrulewidth{0pt}
\lhead{}\chead{}\rhead{}
\rfoot{\vspace*{-2\baselineskip}}


\renewcommand{\vec}{\overrightarrow}

%%%%%%%%%%%%%%%%%%%%%%55\Large

\everymath{\displaystyle}
\begin{minipage}{20cm}
\begin{bclogo}[couleur=brown!40,
  arrondi=0,ombre=true, couleurOmbre=gray!50,
logo =\bcplume,
  noborder=true]
{\textrm{{\Large{Test 1.2.} \Huge{}:
%
%
%%%%%%%%%%%%%%%%%%%%%%%%%%%%%%%%%%%%%%%%%
%%%%%%%%%%%%%%%%%%%%%%%%%%%%%%%%%%%%%%%%%
% TUKAJ VNESI VSEBINO TESTA
\large{Polinomi in racionalna funkcija} \hfill $G-3$}}}
%%%%%%%%%%%%%%%%%%%%%%%%%%%%%
\end{bclogo}
\end{minipage}
\begin{minipage}{20cm}
\begin{bclogo}[couleur=brown!40,ombre=true, couleurOmbre=gray!50,
  arrondi=0,
logo =\bcplume,
  noborder=true]
 \setlength{\parskip}{}{}
{\textsc{{\large{Ime in priimek:}
%%%%%%%%%%%%%%%%%%%%%%%%%%%%%%%%%%%%%%%%%
% TUKAJ VNESI VSEBINO TESTA
\vspace{0.2cm} \rule{15cm}{0.2mm}\hfill }}}
%%%%%%%%%%%%%%%%%%%%%%%%%%%%%%%%%%%%%%%%%
\end{bclogo}
\end{minipage}


\def\tocke2{3+2+3}
\begin{naloga}[\hfill \quad $\tocke2$
\qquad \rightsquigarrow \vert a.\qquad|b.\qquad|c.\qquad|] \addtocounter{tocke}{\tocke2}
Na sliki je polinom tretje stopnje.

a) Zapiši ničle in njihovo stopnjo.

b) Določi začetno vrednost in $p(-1)$.

c) Zapiši predpis za polinom.

\end{naloga}
\definecolor{qqwuqq}{rgb}{0.,0.39215686274509803,0.}
\begin{tikzpicture}[line cap=round,line join=round,>=triangle 45,x=1.0cm,y=1.0cm]
\begin{axis}[
x=1.0cm,y=1.0cm,
axis lines=middle,
ymajorgrids=true,
xmajorgrids=true,
xmin=-4.6945454545454535,
xmax=3.814545454545456,
ymin=-5.3199999999999985,
ymax=2.407272727272729,
xtick={-4.0,-3.0,...,3.0},
ytick={-5.0,-4.0,...,2.0},]
\clip(-4.6945454545454535,-5.32) rectangle (3.814545454545456,2.407272727272729);
\draw[line width=1.pt,color=qqwuqq,smooth,samples=100,domain=-4.6945454545454535:3.814545454545456] plot(\x,{((\x)-1)*((\x)+2)^(2)});
\begin{scriptsize}
\draw[color=qqwuqq] (-3.385454545454544,-8.783636363636361) node {$f$};
\end{scriptsize}
\end{axis}
\end{tikzpicture}





\newpage
\def\tocke18{5}
\begin{naloga}[\hfill \quad $\tocke18$
\qquad \rightsquigarrow \vert a.\qquad\vert b.\qquad|
c.\qquad|] \addtocounter{tocke}{\tocke18}

Zapiši količnik in ostanek pri deljenju polinoma $p(x)=3x^3+4x^2-7x+1$ s polinomom $q(x)=x^2-2x-1.$

\end{naloga}

\vspace{10cm}
\def\tocke6{5}
\begin{naloga}[\hfill \quad $\tocke6$
\qquad \rightsquigarrow \vert a.\qquad| b.\qquad| c.\qquad| ] \addtocounter{tocke}{\tocke6}
Določi ničle polinoma $p(x)=x^4+5x^3+5x^2-5x-6$ s Hornerjem, tako da najprej zapišeš vse kandidate za celoštevilske ničle.
\end{naloga}

\newpage
\def\tocke3{3+3}
\begin{naloga}[\hfill \quad $\tocke3$
\qquad \rightsquigarrow \vert a.\qquad\vert b.\qquad] \addtocounter{tocke}{\tocke3}
S Hornerjevim algoritmom izračunaj:

a)  $p(-2)$ za $p(x)=\dfrac{1}{2}x^3-4x^2+5x-2,$

b) ali je $x^3-2x-1$ deljiv z $x+1.$
 

\end{naloga}

\vspace{10cm}

\def\tocke1{4}
\begin{naloga}[\hfill \quad $\tocke1$
\qquad \rightsquigarrow \vert a.\qquad\vert b.\qquad|
c.\qquad|d.\qquad|] \addtocounter{tocke}{\tocke1}
Reši enačbo: $x^3-4x^2+4x=0$
\end{naloga}
\newpage
\def\tocke6{8}
\begin{naloga}[\hfill \quad $\tocke6$
\qquad \rightsquigarrow \vert a.\qquad| b.\qquad| c.\qquad| ] \addtocounter{tocke}{\tocke6}
Določi ničlo, pol in asimptoto funkcije $f(x)=\dfrac{3x+6}{x-3}$ in jo nariši.
\end{naloga}


\vfill
\mdtheorem[style=theoremstyle1]{kriterij}{Kriterij ocenjevanja}
\begin{kriterij*}[\hfill \v{s}tevilo vseh to\v{c}k na testu: $\arabic{tocke}$]
 $$
 \begin{array}{||c|c|c|c|c|c||
>{}c|c|}
 \hline\hline
 \textup{ocena}& 1&2&3&4&5&\textrm{uspe\v{s}nost v } \%& \textrm{\bf OCENA}\\
  \hline\hline
 \% & [0,45) &[45,60)&[60,75
 ) &[75,90)    &    [90,100]& \mbox{\phantom{ \huge{$\frac{5}{445}$
 15}}
 } 
 & \\
 \hline
 \end{array} \hspace*{0.5cm}\vspace*{-0.3cm} \qrcode[hyperlink,height=0.8in]{http://www.matej.info}
$$
\vspace*{0.1cm}
\end{kriterij*}

\end{document}aloga}[\hfill \quad $\tocke1$
\qquad \rightsquigarrow \vert a.\qquad\vert b.\qquad|
c.\qquad|d.\qquad|] \addtocounter{tocke}{\tocke1}


\end{naloga}


\vfill
\mdtheorem[style=theoremstyle1]{kriterij}{Kriterij ocenjevanja}
\begin{kriterij*}[\hfill \v{s}tevilo vseh to\v{c}k na testu: $\arabic{tocke}$]
 $$
 \begin{array}{||c|c|c|c|c|c||
>{}c|c|}
 \hline\hline
 \textup{ocena}& 1&2&3&4&5&\textrm{uspe\v{s}nost v } \%& \textrm{\bf OCENA}\\
  \hline\hline
 \% & [0,45) &[45,60)&[60,75
 ) &[75,90)    &    [90,100]& \mbox{\phantom{ \huge{$\frac{5}{445}$
 15}}
 } 
 & \\
 \hline
 \end{array} \hspace*{0.5cm}\vspace*{-0.3cm} \qrcode[hyperlink,height=0.8in]{http://www.matej.info}
$$
\vspace*{0.1cm}
\end{kriterij*}

\end{document}